% 预热加余弦衰减与WSD日程的学习率曲线（解析式绘制，参数为教学示例）。
% 余弦：前5%线性预热，之后按余弦降到0.1；WSD：前5%预热，恒定到80%，最后20%线性降到0；
% 虚线为从40%与60%处恒定段检查点分出的短衰减段，长度取各自已训练步数的20%。
\begin{tikzpicture}[figure picture]
  \begin{axis}[
    width=\linewidth,
    height=54mm,
    xmin=0,xmax=100,
    ymin=0,ymax=1.2,
    xlabel={训练进度（占计划总步数的百分比）},
    ylabel={$\eta_t/\eta_{\max}$},
    xtick={0,20,40,60,80,100},
    ytick={0,0.5,1},
    grid=major,
    legend style={at={(0.5,1.03)},anchor=south,legend columns=3,
      /tikz/every even column/.append style={column sep=4mm}},
  ]
    \addplot[FigureBlue,line width=1.1pt,domain=0:5,samples=2,forget plot] {x/5};
    \addplot[FigureBlue,line width=1.1pt,domain=5:100,samples=120]
      {0.1+0.45*(1+cos(180*(x-5)/95))};
    \addlegendentry{预热加余弦}
    \addplot[FigureRed,line width=1.1pt,domain=0:5,samples=2,forget plot] {x/5};
    \addplot[FigureRed,line width=1.1pt,domain=5:80,samples=2,forget plot] {1};
    \addplot[FigureRed,line width=1.1pt,domain=80:100,samples=2] {(100-x)/20};
    \addlegendentry{WSD}
    \addplot[FigureYellow!80!black,line width=1pt,dashed,domain=40:48,samples=2] {(48-x)/8};
    \addlegendentry{分支衰减}
    \addplot[FigureYellow!80!black,line width=1pt,dashed,domain=60:72,samples=2,forget plot] {(72-x)/12};
    \addplot[only marks,mark=*,mark size=1.4pt,FigureStroke,forget plot] coordinates {(40,1) (60,1)};
    \node[anchor=south,font=\figurefont\footnotesize,text=FigureStroke] at (axis cs:50,1.01) {分支检查点};
  \end{axis}
\end{tikzpicture}
